El cuantizador uniforme consiste en una serie de $n\left(b\right)$ niveles equiespaciados de paso $s(\alpha, b) = \frac{\alpha}{n\left(b\right)}$, donde $n\left(b\right) = 2^{b-1}$, \
siendo $b$ el número de bits utilizados para la cuantización. Con cada nivel teniendo el valor:

\begin{equation}
    \label{eq:uniform_quantizer_level_value}
    v_{l_i}\left(\alpha\right) = -\alpha + i \cdot s\left(\alpha, b\right) \text{, con } i = 0, 1, \ldots, n\left(b\right) - 1
\end{equation}

La función de cuantización uniforme puede escribirse como en la ecuación \refeq{uniform_quantizer_signed} y se ve como en la \reffig{uniform_signed}.
\begin{equation}
    \label{eq:uniform_quantizer_signed}
    q(x; \alpha) =
    \begin{cases}
        - \alpha & \text{si $x \leq -\alpha$} \\
        \frac{1}{k} \cdot \lfloor x \cdot k \rfloor & \text{si $-\alpha < x < \alpha$, con $k = \frac{n}{\alpha}$} \\
        v_{l_{n\left(b\right) - 1}}\left(\alpha\right) & \text{si $x \geq \alpha$}
    \end{cases}
\end{equation}

\defaultfigure{quantizers/uniform/q.png}{Quantización uniforme signada.}{uniform_signed}

Para determinar las derivadas respecto a la entrda $\frac{\partial q}{\partial x}$, se puede utilizar el STE (Straight Through Estimator) que consiste en aproximar la función de cuantización como en la \refeq{uniform_quantizer_ste_signed}, como se muestra en la \reffig{uniform_signed_ste}.

\begin{equation}
    \label{eq:uniform_quantizer_ste_signed}
    q\left(x; \alpha\right) =
    \begin{cases}
        - \alpha & \text{si $x \leq -\alpha$} \\
        x \\
        v_{l_{n\left(b\right) - 1}}\left(\alpha\right) & \text{si $x \geq \alpha$}
    \end{cases}
\end{equation}

\defaultfigure{quantizers/uniform/q_ste.png}{Aproximación de la cuantización uniforme signada con STE.}{uniform_signed_ste}

Con lo cual resulta evidente que la derivada de la función de cuantización uniforme es la derivada de la función identidad, es:

\begin{equation}
    \label{eq:uniform_quantizer_ste_signed_derivative}
    \frac{\partial q}{\partial x}\left(x; \alpha\right) =
    \begin{cases}
        0 & \text{si $x \leq -\alpha$} \\
        1 & \text{si $-\alpha < x < \alpha$} \\
        0 & \text{si $x \geq \alpha$}
    \end{cases}
\end{equation}

Ahora queda resolver la derivada con respecto a $\alpha$ de la función de cuantización uniforme, $\frac{\partial q}{\partial \alpha}$.

Si observamos la \reffig{uniform_signed}, podemos ver que la función de cuantización uniforme signada se puede expresar como:

\begin{equation}
    \label{eq:uniform_quantizer_signed_dq_dalpha}
    \frac{\partial q}{\partial \alpha}\left(x; \alpha\right) = \frac{q(x; \alpha)}{\alpha}
\end{equation}

Pues para un $x$ dado, si se incrementa $\alpha$, la función de cuantización se incrementa en la misma proporción.

